\scanpage{76}
% archive foliation: 75

Il est immédiat que \deleted{\unsure{cet}}\deleted{\unsure{iso}}
l'isom ne dépend pas du choix des plongements $L\hookrightarrow E$,
$M\hookrightarrow E$.

\uline{Proposition} La contrainte de commutativité précédente ne dépend
pas du choix de $X\in\operatorname{Ob}C$.

\deleted{Considérons en effet un objet $X'$ de}
[
Il est clair que cela ne \unsure{change pas quand on}
remplace $X$ par un objet isomorphe
[\unsure{il faudrait} une petite démonstration : dans une gr-catégorie
$B$, les \deleted{\unsure{foncteurs}} autoisomorphismes
\deleted{\unsure{qui}} transportent une
\deleted{\unsure{donnée}}\inserted{contrainte}
de comm. (\unsure{stricte ou non}) en elle-même \ldots
]]
\margin{$\unsure{\pi_0}$}

\deleted{On peut \unsure{donc toujours supposer}}
Soit $X'$ un autre choix, défini par un $P\in\mathbb P(O/A)$.
On peut supposer $P\subset E$, soit $O'=\operatorname{End}_O(P)$,
de sorte que $O'$ est maintenant plongé dans $E$
(c'est l'ens des \deleted{\unsure{Aut}} $\lambda\in E$ tels que
$\lambda P\subset P$) Ainsi $P$ est un \deleted{\unsure{biidé}}
$(O',O)$-biidéal. On a un diagramme comm.

\diagram{
\begin{tikzcd}[column sep=large,row sep=large]
& \uline{\operatorname{Eq}}(C,C)
\arrow[dl,"\alpha_X"']\arrow[dr] &\\
\begin{gathered}\text{Catégorie des}\\(O,O)\text{ biidéaux}\\\text{spéciaux}\end{gathered}
\arrow[rr,"L\mapsto P\deleted{\unsure{\otimes}}LP^{-1}"]
&&\begin{gathered}\text{Catégorie des}\\(O',O')\text{-biidéaux}\\\text{spéciaux}\end{gathered}
\end{tikzcd}
}

Il suffit donc de prouver que \unsure{le fonct. $P.L$}
\deleted{\unsure{est compatible}}\inserted{\unsure{commute}}
\deleted{\unsure{aux}}\inserted{avec les}
contraintes de commutativité
\diagram{
\begin{array}{ccc}
LM&&\\[-2pt]
\scriptstyle c_X=1\quad\Vert\ \deleted{\unsure{\simeq}}
&\text{transportées}&
\begin{array}{ccc}
P(LM)P^{-1}&=&(PLP^{-1})(PMP^{-1})\\
\Vert\ \scriptstyle\unsure{F}(c_X)=1&&\Vert\ \scriptstyle c_{X'}=1\\
P(ML)P^{-1}&=&(\deleted{\unsure{P}}PMP^{-1})(PLP^{-1})
\end{array}\\[-2pt]
ML&&
\end{array}
}
OK.
